
\documentclass[../../../physics-standard_questions_all.tex]{subfiles}


\begin{document}

図のように，電池E（起電力$V_0=5.0$\,V），抵抗R，コンデンサC，
およびスイッチSからなる回路があり，Rの電流を$I$，Cの電荷を$Q$と表す．
はじめSは開いており，Cは帯電していない（$Q=0$）．
時刻$t=0$にSを閉じる．

\begin{enumerate}

    \item $t=0$直後の$I$の値は$I_0=2.0$\,mAであった．
    また，充分に時間が経過した後，$Q$の値は$Q_\infty=15$\,\textmu Cで一定となった．
    このことから，Rの抵抗値$R$，およびCの容量$C$を，それぞれ数値で求めよ．

\end{enumerate}

以下は文字式で答えよ．

\begin{enumerate}[resume]

    \item $Q=kQ_\infty$\,($0\leqq k \leqq 1$)となった瞬間の$I$の値を，
    $k$を用いて表せ．

    \item $I$，$Q$を，それぞれ$t$の関数として表せ．      
\end{enumerate}

\vspace{0.2\baselineskip}
{\centering
    \includegraphics{\subfix{fig/fig_em_cb_04/fig_em_cb_04_00.pdf}}
\par}
\vspace{0.2\baselineskip}

\end{document}



